% ============================================================
% CAPÍTULO 8: NUCLEÓTIDOS Y ÁCIDOS NUCLEICOS
% ============================================================
% CONTENIDO BASADO EN: Unidad de Aprendizaje II, Tema "Nucleótidos y ácidos nucleicos"
% PROPÓSITO: Reconocer las estructuras de los nucleótidos y ácidos nucleicos, y señalar sus funciones.
% ============================================================

\section{Estructura de los nucleótidos}

% --- CONSEJO: Define el nucleótido como la unidad básica de los ácidos nucleicos. ---
\resumebox{definicion}{Nucleótido}{Un nucleótido es una molécula compuesta por tres partes: una base nitrogenada, un azúcar (pentosa) y uno o más grupos fosfato.}

\subsection{Componentes de un nucleótido}
% --- CONSEJO: Explica cada componente por separado. ---
\begin{enumerate}
    \item \textbf{Base nitrogenada:} Molécula heterocíclica que contiene nitrógeno. Se dividen en dos grupos:
    \begin{itemize}
        \item \textbf{Purinas:} Adenina (A) y Guanina (G).
        \item \textbf{Pirimidinas:} Citosina (C), Timina (T) y Uracilo (U).
    \end{itemize}
    \item \textbf{Pentosa (azúcar de 5 carbonos):}
    \begin{itemize}
        \item \textbf{Ribosa:} Presente en el ARN.
        \item \textbf{Desoxirribosa:} Presente en el ADN. Difiere de la ribosa en la ausencia de un grupo hidroxilo (-OH) en el carbono 2'.
    \end{itemize}
    \item \textbf{Grupo fosfato:} Uno o más grupos de ácido fosfórico unidos al azúcar.
\end{enumerate}

% --- CONSEJO: Puedes añadir una figura con la estructura general de un nucleótido y otra con las bases nitrogenadas. ---

\section{Estructura de los ácidos nucleicos}

% --- CONSEJO: Explica cómo se unen los nucleótidos para formar polímeros (ácidos nucleicos) a través de enlaces fosfodiéster. ---
\resumebox{definicion}{Ácido nucleico}{Es un polímero de nucleótidos unidos por enlaces fosfodiéster entre el grupo fosfato de un nucleótido y el grupo hidroxilo del azúcar del siguiente.}

\subsection{Ácido desoxirribonucleico (ADN)}
% --- CONSEJO: Describe la estructura de doble hélice del ADN, propuesta por Watson y Crick. ---
\begin{itemize}
    \item Es una doble hélice formada por dos cadenas polinucleotídicas antiparalelas.
    \item Las bases nitrogenadas se aparean por puentes de hidrógeno (A=T, G$\equiv$C).
    \item La información genética está almacenada en la secuencia de bases.
\end{itemize}

\subsection{Ácido ribonucleico (ARN)}
% --- CONSEJO: Describe las diferencias con el ADN y los diferentes tipos de ARN. ---
\begin{itemize}
    \item Generalmente es una cadena sencilla.
    \item La base nitrogenada es A, G, C y U (uracilo) en lugar de timina.
    \item El azúcar es ribosa.
    \item Tipos principales y sus funciones:
    \begin{itemize}
        \item \textbf{ARN mensajero (ARNm):} Lleva la información del ADN a los ribosomas para la síntesis de proteínas.
        \item \textbf{ARN de transferencia (ARNt):} Transporta aminoácidos a los ribosomas durante la traducción.
        \item \textbf{ARN ribosómico (ARNr):} Es un componente estructural y catalítico de los ribosomas.
    \end{itemize}
\end{itemize}

\section{Funciones de los ácidos nucleicos}

% --- CONSEJO: El programa pide señalar las funciones. ---
\begin{itemize}
    \item \textbf{Almacenamiento de información genética:} El ADN almacena toda la información necesaria para la vida de un organismo.
    \item \textbf{Transmisión de la información genética:} El ADN se replica y se transmite de padres a hijos.
    \item \textbf{Expresión de la información genética:} El ARN permite la síntesis de proteínas a partir de la información contenida en el ADN (dogma central de la biología molecular).
    \item \textbf{Funciones reguladoras y catalíticas:} Algunos ARN (ribozimas) tienen actividad catalítica, y otros tipos de ARN (ARN pequeños) regulan la expresión génica.
    \item \textbf{Metabolismo energético:} Los nucleótidos como el ATP (adenosín trifosfato) y el GTP (guanosín trifosfato) son moléculas portadoras de energía celular.
    \item \textbf{Señalización celular:} Los nucleótidos cíclicos como el AMPc (adenosín monofosfato cíclico) son mensajeros secundarios en rutas de señalización.
\end{itemize}

\section{Importancia biotecnológica de los ácidos nucleicos}

% --- CONSEJO: Esta es una sección fundamental para conectar con la biotecnología. ---
\begin{itemize}
    \item \textbf{Ingeniería genética:} El conocimiento de la estructura y función de los ácidos nucleicos es la base de la ingeniería genética. Permite manipular el ADN para:
    \begin{itemize}
        \item Crear organismos genéticamente modificados (OGM).
        \item Producir proteínas recombinantes (insulina, hormonas, enzimas).
    \end{itemize}
    \item \textbf{Diagnóstico molecular:} Técnicas como la PCR y la secuenciación de ADN permiten detectar enfermedades genéticas, agentes patógenos y mutaciones.
    \item \textbf{Biotecnología médica:} Desarrollo de terapias génicas y ARN de interferencia (ARNi) para tratar enfermedades.
    \item \textbf{Biotecnología forense y de identificación:} Uso de perfiles de ADN (huellas genéticas) en investigaciones criminales y pruebas de paternidad.
\end{itemize}